% !TeX TS-program = xelatex
% 虚构演示简历 — 仅用于 hellojobs 测试数据
\documentclass{resume-photo}
\usepackage{xeCJK}
\setCJKmainfont{Noto Serif CJK SC}
\usepackage{fontspec}
\usepackage{enumitem}
\setlist[itemize]{itemsep=0pt, topsep=1pt, parsep=0pt, partopsep=0pt}

\ResumeName{陆凯}

\begin{document}

\ResumeContacts{
  138-0028-0028,%
  \ResumeUrl{mailto:lukai@cqu.edu.cn}{lukai@cqu.edu.cn},%
  \textnormal{重庆大学 | 物联网工程 · 学士 | 2020—2024}%
}

\ResumeTitle

\noindent 边缘计算与 IoT 平台，熟悉 MQTT、KubeEdge 与设备孪生；园区路灯边缘网关规则引擎 CPU 占用下降 19\%。注重可观测性与文档沉淀，习惯在评审阶段拆解风险；参与线上复盘与跨团队联调，英语 CET-6，可阅读官方文档与 RFC（演示数据）。

\section{教育背景}
\ResumeItem
{重庆大学}
[\textnormal{微电子与通信工程学院 | 物联网工程 | 学士}]
[2020.09—2024.06]

\begin{itemize}
  \item 嵌入式与通信协议课程。
\end{itemize}


\ResumeItem
{学术交流与在线课程}
[\textnormal{暑校 / MOOC / 厂商认证（演示）}]
[2022—2025]

\begin{itemize}
  \item 完成《深入理解计算机系统》配套实验与 MIT 6.828 / 6.S081 公开课部分章节跟做。
  \item Coursera / edX 分布式系统、云计算相关专项课程结业证书 4 门（虚构编号）。
  \item 通过 AWS SAA / 阿里云 ACP / Redis 开发者等认证中的 1--2 项（简历测试数据）。
  \item 参加 CCF CCSP / 华为 ICT / 百度之星等学科竞赛集训营并完成全部作业。
\end{itemize}



\section{专业技能}
\begin{itemize}
  \item 熟悉 Go、MQTT、EMQX 规则引擎与 TimescaleDB。
  \item 掌握 Docker、KubeEdge 云边协同部署。
  \item 了解 Modbus/Bacnet 与 OPC UA 网关桥接。
  \item 熟悉 Linux 日常运维与 Shell，具备日志检索与 CPU/内存突增初步定位能力。
  \item 掌握 Git / CI 与 Code Review，了解 Docker 与 Kubernetes 基本编排与滚动发布。
  \item 了解 OAuth2 / JWT、接口幂等与 REST 规范，具备单元测试与集成测试习惯（JUnit/pytest 等）。
  \item 能阅读英文技术文档，具备跨团队沟通与需求澄清能力（测试简历）。
\end{itemize}

\section{实习经历}
\ResumeItem
{海尔智家}
[\textnormal{IoT 平台实习生}]
[2024.03—2024.08]

\begin{itemize}
  \item \textbf{职责}: 设备影子同步延迟监控与告警。
  \item \textbf{成果}: P99 延迟从 1.4s 降至 0.6s（华东区抽样）。
\end{itemize}


\ResumeItem
{拼多多 | 交易平台}
[\textnormal{后端实习生}]
[2024.03—2024.08]

\begin{itemize}
  \item \textbf{活动}: 参与大促「万人团」库存预占与释放逻辑开发与压测。
  \item \textbf{一致性}: 使用分布式锁 + 版本号控制超卖风险，压测零超卖。
  \item \textbf{观测}: 搭建 Grafana 看板跟踪关键指标，值班期间 0 P1 事故。
  \item \textbf{复盘}: 参与事故复盘 2 次，输出改进项并跟踪闭环。
  \item \textbf{代码}: MR 平均评论轮次 4 轮，注重可读性与边界条件注释。
\end{itemize}



\section{项目经历}
\ResumeItem
{边缘规则热更新}
[\textnormal{毕业设计}]
[2023.10—2024.05]

\begin{itemize}
  \item 基于 WebAssembly 沙箱执行用户脚本，崩溃隔离率 100\%。
\end{itemize}


\ResumeItem
{低代码表单引擎}
[\textnormal{拖拽 + JSON Schema（演示）}]
[2023.07—2023.12]

\begin{itemize}
  \item \textbf{前端}: 基于 React 实现组件注册表与撤销重做栈。
  \item \textbf{后端}: 动态表单元数据存储与版本迁移脚本。
  \item \textbf{生态}: 导出 JSON 与后端校验规则一致，减少联调返工。
  \item \textbf{采用}: 在内 3 个业务线试点，表单搭建时间缩短 70\%（演示）。
\end{itemize}



\section{个人荣誉}
\begin{itemize}
  \item 挑战杯 重庆市特等奖
  \item 重庆大学 优秀本科毕业生
  \item 全国计算机等级考试四级（网络工程师）（演示）
  \item 校级「优秀学生干部 / 优秀团员」或技术社团骨干（综合测评前 15\%）
  \item 开源贡献：向知名项目提交被合并的 PR（文档/测试/feature）（演示）
\end{itemize}

\end{document}
